\begin{python0}
from solutions import *; clear()
\end{python0}
\sectionthree{Example: Cellular automata}

A \textbf{cellular automata (CA)} is simply a grid of values where the value at a point in the grid can change its value. The way such a value \verb!v! changes depends on the values near \verb!v!.

CAs are studied in math, CS, physics, biology, social science, etc. You name it. They can be as practical as image processing where they are used to remove noise from images to create cleaner images. Yet they can be as abstract and complex as you like -- they appear in AI, dynamical systems, and chaos theory. You can find lots of information about cellular automata on the web -- go ahead and check out the CA entry at wikipedia.

Let me be more specific by looking at a simple example. Suppose we look at a 1-dimensional CA with these values:
\begin{center}
0,0,1,1,0,1,1,1
\end{center}
This CA is made up of 8 cells. The values are either 0 or 1. So a CA can be as simple as an array.

Now suppose the value at a cell changes according to these rules:

\begin{itemize}
\item
  If the value is 1 and together with the values on its left and right,
  there are three or one 1s, then the value becomes 0. Otherwise it
  stays as 1. In other words, if the value is 1 and either the left or
  right neighbor is 1 (but not both), then 1 stays as 1. Otherwise it
  becomes 0. You can think of it this way:

  \begin{itemize}
  \item
    Companionship: If 1 has exactly one companion, he/she/it lives on.
    If 1 has no companion, it dies.
  \item
    Overcrowding: If 1 has too many companions, overcrowding kills
    he/she/it.
  \end{itemize}
\item
  If this value is 0, and together with the values on its left and
  right, there are two 1s, then the value becomes 1. Otherwise it stays
  as 0. You can think of it this way:

  \begin{itemize}
  \item
    Reproduction: If a spot is available, then the 1 on the left and
    right produces a 1. A reproduction (i.e., $0 \rightarrow 1$) occurs only when
    there are two adjacent 1s next to the 0.
  \end{itemize}
\end{itemize}

For instance, the overcrowding rule gives is this:

\begin{center}
1,1,1
$\downarrow$
0
\end{center}

So applying this rule to the value at index 6 we get

\begin{consolethree}[escapeinside=||]
0,0,1,1,0,1,1,1
            |$\downarrow$|
0,0,1,1,0,1,0,1
\end{consolethree}

I'll write the rule as $111 \rightarrow 0$ instead of

\begin{consolethree}[escapeinside=||]
1,1,1
  |$\downarrow$|
  0
\end{consolethree}

Here are more examples:

\begin{itemize}
\item
  Using the rule $001 \rightarrow 0$, we get

\begin{consolethree}[escapeinside=||]
0,0,1,1,0,1,1,1
  |$\downarrow$|
_,0,_,_,_,_,_,_
\end{consolethree}

\item
  Using the rule $011 \rightarrow 1$, we get

\begin{consolethree}[escapeinside=||]
0,0,1,1,0,1,1,1
    |$\downarrow$|
_,_,1,_,_,_,_,_
\end{consolethree}

\item
  Using the rule $101 \rightarrow 1$, we get

\begin{consolethree}[escapeinside=||]
0,0,1,1,0,1,1,1
        |$\downarrow$|
_,_,_,_,1,_,_,_
\end{consolethree}
\end{itemize}

In general the new i-th value depends on the current (i-1)-th, i-th, (i+1)th values:

\begin{consolethree}[escapeinside=||]
_,_,_,_,?,_,_,_
      |$\downarrow$|  |$\downarrow$|  |$\downarrow$|
      |$\searrow$| |$\downarrow$| |$\swarrow$|
                 |$\downarrow$|
_,_,_,_,?,_,_,_
\end{consolethree}
The value at the ends (the first and last) do not have two neighbors. So xwe will not change the values at the left and right end points:

\begin{consolethree}[escapeinside=||]
0,0,1,1,0,1,1,1
|$\downarrow$|             |$\downarrow$|
0,_,_,_,_,_,_,1
\end{consolethree}

(There are other ways to compute the value of the cells at end points. For instance you can view the CA as being wrapped around at the end points so that the new value for the first cell depends on the values of first, second, and last cell).

So using the above rules we get the following behavior
\begin{consolethree}[escapeinside=||]
0,0,1,1,0,1,1,1

0,0,1,1,1,1,0,1
\end{consolethree}
You can think of the above as an evolving CA that changes with time. With three times we get

\begin{center}
0,0,1,1,0,1,1,1 (time 0)

0,0,1,1,1,1,0,1 (time 1)

0,0,1,0,0,1,1,1 (time 2)
\end{center}

etc.

We will use CA to generate some ASCII art by printing the values of the cells in 1D CA: If the value is 0, we print a space and if the value is 1 we print X. We will let our CA run through a certain number of time steps, printing the CA for each time step. For instance the above CA goes through three time steps:
\begin{center}
0,0,1,1,0,1,1,1

0,0,1,1,1,1,0,1

0,0,1,0,0,1,1,1
\end{center}

and we print

\begin{consolethree}[escapeinside=||]
+--------+
|  XX XXX|
|XXXX  X |
| X  XXX |
+--------+
\end{consolethree}

We will use a 1D CA of size 2 * n + 1 where n is an input from the user. We will run this for n time steps. We will start off with a CA with a 1 in the middle of the 1D array and 0 elsewhere. The set of rules to use is:

\begin{center}
\begin{tabular}{rcl}
000 & $\rightarrow$ & 0 \\
001 & $\rightarrow$ & 1 \\
010 & $\rightarrow$ & 0 \\
011 & $\rightarrow$ & 1 \\
100 & $\rightarrow$ & 1 \\
101 & $\rightarrow$ & 0 \\
110 & $\rightarrow$ & 1 \\
111 & $\rightarrow$ & 0 \\
\end{tabular}
\end{center}

Here's the pseudocode:

\begin{tabular}{l}
declare an array ca of size 2*500 + 1.\\
note that the maximum size of the ca is 2*500 + 1.\\
declare an array t of size 2*500 + 1.\\

get n from user (at most 500), size of the ca is 2 * n + 1.\\
note that we will only use ca[0], \ldots, ca[2*n].\\
set the values in ca to all 0s except for a 1 in the middle.\\

// time = 0\\
print ca (for value 1 print 'X' and for value 0 print ' '; use a loop).\\
for time = 1, 2, ..., n - 1:\\
\{\\
\tab[3em]{for each value in ca,}\\
\tab[3em]{apply above rules and fill corresponding value in t}\\
\tab[3em]{copy values in t back to ca (use a loop).}\\
\tab[3em]{print ca (for value 1 print 'X' and for value 0 print ' ').}\\
\}\\
\end{tabular}

Note that you can specify n = 500 for a maximum CA of size 2 * 500 + 1 = 1001. However, you can't see the ASCII art clearly on your console window because of wraparound. You can probably right-click and choose a smaller font size and larger window size and then you can see the ASCII art up to about n = 100. For larger sizes, you can save the output to a document (say MS Word), choose a really tiny font and a larger page size and print it out.

NOTE: Recall that a number such as 1425 is just 1*1000 + 4*100 + 2*10 + 5. You can extract the 1, 4, 2, and 5 from 1425 by using integer division / and integer mod \%. This is viewing an integer ``written in base 10''. If I give you 1,0,1 or any sequence of zeroes and ones, then it's sometimes convenient to convert that into a single integer. For instance you can convert 1,0,1 into 5 using ``binary representation'' -- I'll explain the conversion in a bit. But first, why would you want to do that? Well,

compare this

\begin{consolethree}[escapeinside=||]
int a = 1, b = 0, c = 1; // 1,0,1

if (a == 1 && b == 0 && c == 1)
{
    ...
}
\end{consolethree}

with this:

\begin{consolethree}[escapeinside=||]
int d = 5; // 1,0,1 converted to 5

if (d == 5)
{
    ...
}
\end{consolethree}

Clearly the second code fragment is simpler because after the conversion there's only one variable and not three. Also, because you are comparing a single integer value, you can use a switch:

\begin{consolethree}[escapeinside=||]
switch (d)
{
    case 5:
        ...
        break;
}
\end{consolethree}

But what is a nice way to convert 1,0,1 to an integer? You use the same idea as base 10 representation of numbers but instead of powers of 10, you use powers of 2:
\begin{center}

\[1,0,1 \rightarrow 1 \times 4 + 0 \times 2 + 1 \times 1\]
\end{center}

Here's another example:
\begin{center}

\[10111 \rightarrow 1 \times 16 + 0 \times 8 + 1 \times 4 + 1 \times 2 + 1 \times 1\]
\end{center}

In the case of our CA, the three ``bits'' determining how to change a bit can be converted to an integer. The rules:

\begin{center}
\begin{tabular}{rcl}
000 & $\rightarrow$ & 0 \\
001 & $\rightarrow$ & 1 \\
010 & $\rightarrow$ & 0 \\
011 & $\rightarrow$ & 1 \\
100 & $\rightarrow$ & 1 \\
101 & $\rightarrow$ & 0 \\
110 & $\rightarrow$ & 1 \\
111 & $\rightarrow$ & 0 \\
\end{tabular}
\end{center}

if you convert the 0s and 1s on the left to base 10 numbers, becomes

\begin{center}
\begin{tabular}{rcl}
0 & $\rightarrow$ & 0 \\
1 & $\rightarrow$ & 1 \\
2 & $\rightarrow$ & 0 \\
3 & $\rightarrow$ & 1 \\
4 & $\rightarrow$ & 1 \\
5 & $\rightarrow$ & 0 \\
6 & $\rightarrow$ & 1 \\
7 & $\rightarrow$ & 0 \\
\end{tabular}
\end{center}

